\documentclass[12pt,a4paper]{article}
%%%%%%%%%%%%%%%%%%%%%%%%%%%%%%%%%%%%%%%%%%%%%%%%%%%%%
%\usepackage[utf8]{inputenc}
\usepackage[portuguese]{babel}
\addto{\captionsportuguese}{%
  \renewcommand{\refname}{\hfill \large REFERÊNCIAS \hfill} % V V V
}
\addto{\captionsportuguese}{% 
  \renewcommand{\contentsname}{\hfill SUMÁRIO \hfill}} %\centering não tava funcionando sabe deus por que, usei o \hfill pra enjambrar
    
\usepackage[T1]{fontenc}
\usepackage{fontspec}
\usepackage{amsmath,esint}
\usepackage{amsfonts}
\usepackage{amssymb}
\usepackage{makeidx}
\usepackage{graphicx}
\usepackage{lettrine}
\usepackage{lmodern}
\usepackage{dirtytalk}
\usepackage{setspace}
    \linespread{1.5}
\usepackage{indentfirst}
    \setlength{\parindent}{1cm}
    \setlength{\parskip}{0pt}
\usepackage{pdfpages}
\usepackage[left=3cm,right=2cm,top=3cm,bottom=2cm]{geometry}
\usepackage{hyperref}
\usepackage{float}
\usepackage{multirow}
\usepackage{booktabs}
\usepackage{subfigure}
\usepackage{fancyhdr}
\usepackage{soul} % Pacote que permite grifar palavras no texto

\usepackage[ %faz a citação ser do tipo autor data, faz parecer et al na citação, mas na referencia mostra todos os autores e deixa o titulo em negrito na referencia
	alf,
	abnt-repeated-title-omit=yes,	
	abnt-emphasize=bf,
	abnt-etal-list=0
	]
	{abntex2cite}
\citebrackets()%faz a citação ser entre parenteses e não entre colchetes

%-----------------------------------------------------------
\hypersetup{
%links coloridos
   pdfborder={0 0 0},
   colorlinks = true,  
   linkcolor = black,          % Cor dos links internos
   citecolor = black        % Cor dos links das referências bibliográficas
}
\setmainfont{Times New Roman}
\pagestyle{fancy}
% Clear the header and footer
\fancyhead{}
\fancyfoot{}
\fancyhead[R]{\thepage}
\renewcommand{\headrulewidth}{0pt}%
\setlength{\headheight}{14.49998pt}
%-----------------------------------------------------------

%------Redefinido Tamanho de Fontes-------------------------

\renewcommand{\large}{\fontsize{14}{15}\selectfont}  % large vira 14pt
\usepackage{titlesec}
    \titleformat{\section}
      {\normalfont\fontsize{14}{15}\bfseries}{\thesection}{1em}{}

%-----------------------------------------------------------


\begin{document}

%___________________________CAPA___________________________%
    \begin{titlepage}

    \begin{figure}[H]
        \centering
        \includegraphics[width=0.25\textwidth]{images/brasaooficialcolorido.png}
    \end{figure}

        \begin{center}
                \textbf{PROJETO DE ESTÁGIO SUPERVISIONADO}
            
                \textbf{CURSO DE LICENCIATURA EM FÍSICA} 
        \end{center}
   
    \vspace*{0.2cm}
    \begin{center}

        \textbf{NOME}
       
    \end{center}
    \vspace*{5cm}

    \begin{center}
    \large
         \textbf{PROJETO}
    \end{center}
      
    \vfill  
    
    \begin{center}
        \textbf{RIO DO SUL}

        \textbf{ano}
        
    \end{center}
    
    

    \end{titlepage}



%______________________Folha de Rosto__________________________%

\thispagestyle{empty} %sem numeração de página

    \begin{center}
        \textbf{INSTITUTO FEDERAL CATARINENSE}
        
        \textbf{\textit{CAMPUS} RIO DO SUL}
    \end{center}
    
    \vspace*{0.2cm}
    
    \begin{center}

        \textbf{NOME}
       
    \end{center}

    \vspace{5 cm}

    \begin{center}
        \large
        \textbf{PROJETO DE ESTÁGIO SUPERVISIONADO}
    \end{center}

    \vspace{3 cm}
    
    \hfill\begin{minipage}{\dimexpr\textwidth-8,5 cm}
        \linespread{1,0}\selectfont
        \large 
            Projeto de Trabalho de Curso do Curso de Licenciatura em Física do Instituto Federal Catarinense (IFC) Campus Rio do Sul.\\

            Professor Orientador: XXXX

    \xdef\tpd{\the\prevdepth}
    \end{minipage}

    \vfill

    \begin{center}
        \textbf{RIO DO SUL\\
        ano}
    \end{center}





\newpage

%______________________Folha em Branco__________________________%
\thispagestyle{empty} 

\

\newpage

%______________________________Sumário_________________________%
\thispagestyle{empty} 
\tableofcontents
\newpage
%_____________________________Introdução_______________________%
\section{INTRODUÇÃO}



\newpage

%___________________________Problemas___________________________%
\section{DEFINIÇÃO E DELIMITAÇÃO DO PROBLEMA DE PESQUISA}



\newpage

%_________________________Objetivos____________________________%
\section{PRESSUPOSTOS E OBJETIVOS}



\newpage

%________________________Base Teórica___________________________%
\section{BASE TEÓRICA}



\newpage

%________________Procedimentos Metodológicos____________________%
\section{PROCEDIMENTOS METODOLÓGICOS}

\input{Capítulos/Procedimentos Metodológicos}

\newpage

%________________________Cronograma___________________________%
\section{CRONOGRAMA}



\newpage

\bibliography{bibliografia}  %Preencha o arquivo bibliografia.bib com as referência adequadamente formatadas (pesquise sobre arquivos de extensão .bib) para que aparecem já em formato ABNT

\end{document}
